% ============================================================================
% Fichier  : front/02_mode_emploi.tex (Mode d'emploi)
% Version  : 1.0 -- novembre 2026
% ============================================================================
% 30 h de cours magistral = 15 séances de 2 h ; TD (fiches, exercices) et TP
% (dix labs) faits par l'étudiant chez lui, hors des 30 h.
% ============================================================================

\chapter*{Mode d'emploi du cours}
\addcontentsline{toc}{chapter}{Mode d'emploi du cours}
\markboth{Mode d'emploi}{Mode d'emploi}

\section*{Structure générale}

Le cours compte un prologue et six parties, soit seize chapitres, pour
30~heures de cours magistral réparties en quinze séances de deux heures.
Les travaux dirigés et les travaux pratiques se font à la maison, hors de
ces 30~heures~: pour les travaux dirigés, la fiche de questions de chaque
chapitre et ses exercices de fin de chapitre~; pour les travaux pratiques,
les dix laboratoires de NDAP SUKULU, d'une heure chacun. Les corrigés des
exercices sont réunis dans un livret séparé.

\begin{tableaucours}{cL{6.4cm}cL{3.2cm}}{%
    \tableauHeader{Partie} &
    \tableauHeader{Thème} &
    \tableauHeader{Chap.} &
    \tableauHeader{Niveau}
}
0 & Prologue~: le schéma directeur dort dans un tiroir & P & introductif \\
I & Fondements~: stratégie, alignement, valeur & 1--3 & fondamental \\
II & Le contexte international & 4--5 & fondamental \\
III & Le contexte national & 6--7 & fondamental \\
IV & Méthode~: élaborer le schéma directeur & 8--12 & intermédiaire \\
V & Mise en œuvre~: piloter et maîtriser & 13--14 & intermédiaire \\
VI & Défis et recherche & 15--16 & expert à recherche \\
\end{tableaucours}

\section*{Les quinze séances}

\begin{tableaucours}{cL{6.8cm}L{4.2cm}}{%
    \tableauHeader{Séance} &
    \tableauHeader{Contenu} &
    \tableauHeader{Travail à la maison}
}
1 & Prologue et Chap.~1~: stratégie, SI, numérique & Fiches P et 1 \\
2 & Chap.~2~: les théories de la stratégie & Fiche 2, lab 01 \\
3 & Chap.~3~: l'alignement et la valeur & Fiche 3, lab 02 \\
4 & Chap.~4~: normes et référentiels & Fiche 4 \\
5 & Chap.~5~: doctrines et stratégies numériques & Fiche 5 \\
6 & Chap.~6 et~7~: stratégie numérique du Cameroun, cadre juridique &
Fiches 6 et 7, lab 03 \\
7 & Chap.~8~: cadrer et gouverner la démarche & Fiche 8 \\
8 & Chap.~9~: diagnostiquer l'existant & Fiche 9, lab 04 \\
9 & Chap.~10~: fixer la cible & Fiche 10, lab 05 \\
10 & Chap.~11~: portefeuille et trajectoire & Fiche 11, lab 06 \\
11 & Chap.~12~: chiffrer, financer, acheter, doter & Fiche 12, lab 07 \\
12 & Chap.~13~: piloter l'exécution & Fiche 13, lab 08 \\
13 & Chap.~14~: risques, sécurité, continuité, conformité & Fiche 14,
lab 09 \\
14 & Chap.~15~: données, IA, plateformes, sobriété & Fiche 15 \\
15 & Chap.~16~: auditer un schéma directeur & Fiche 16, lab 10 \\
\end{tableaucours}

Cette répartition est indicative~; l'enseignant l'adapte au calendrier et au
rythme du groupe.

\section*{Le gabarit des chapitres}

Chaque chapitre suit le même plan~: une épigraphe, le niveau et les
prérequis, les objectifs, une situation d'ouverture tirée de l'un des trois
cas, des sections numérotées, ce qu'il faut retenir, des points de
contrôle, trois à quatre exercices gradués, des lectures et une
transition.

\section*{Les encadrés}

\begin{tableaucours}{L{3.6cm}L{8.6cm}}{%
    \tableauHeader{Encadré} &
    \tableauHeader{Rôle}
}
Objectifs & Ce que l'on doit savoir faire à la fin du chapitre \\
Définition & Les notions, en gras dans le texte et indexées \\
Violation typique & Une erreur fréquente et la règle qui l'évite \\
Information & Un complément, un renvoi, une nuance \\
Mission (ANPP, LiZbiZ, BANQUE SANAGA) & Le cas~: les constatations, puis
l'analyse, avec un niveau de confiance \\
Démarche & Les étapes d'une procédure, dans l'ordre \\
Lecture & Un décompte de cotes \OUI{}~/ \NON{}~/ \NEK{} et sa lecture \\
Métier & Ce que la notion change pour chacun des métiers du cours \\
Texte juridique & Un texte cité avec sa référence et sa portée \\
Ce qu'il faut retenir & Les faits et les règles du chapitre \\
Exercice & Trois niveaux de difficulté~: \niveaubase, \niveaumoyen,
\niveauexpert \\
\end{tableaucours}

\section*{Les conventions}

\begin{itemize}
    \item Les cotes~: \OUI{} (fait établi par une pièce), \NON{} (écart
          établi par une pièce), \NEK{} (connaissance insuffisante pour
          trancher). Un \NEK{} est de preuve quand une pièce manque, de
          critère quand la référence manque.
    \item Constatation et analyse sont toujours séparées~; l'analyse porte
          un niveau de confiance (faible, moyenne, élevée).
    \item Les textes~: chaque loi, décret ou règlement est cité avec son
          numéro et sa date~; les normes sous licence sont paraphrasées~;
          la source de chaque texte figure à l'annexe~\ref{ann:textes}.
    \item Les montants sont en francs CFA. Un « milliard » vaut mille
          millions.
    \item Les cas sont fictifs~: noms en majuscules, villes réelles.
\end{itemize}
